% Source: https://artofproblemsolving.com/community/c3h1885113_tricky_number_theory__problems
% Stable identity: aops-topic-1885113-q02
(a) $\dfrac {a}{b}<1$ means $a<b$, so the digits are different and the fraction is determined by the pair $\{a,b\}$ of digits. There are $\dbinom92=\boxed{36}$ such fractions.

(b) Every such fraction reduces to a fraction $\dfrac {p}{q}$ in lowest terms with $1\le p<q\le9$, and every such reduced fraction is itself one of the fractions in (a). So the number of different values equals the number of reduced fractions $\dfrac {p}{q}$ with $2\le q\le9$, $1\le p<q$, $\gcd(p,q)=1$. For $q=2,3,\dots,9$ the numbers of such $p$ are
\[
1,\ 2,\ 2,\ 4,\ 2,\ 6,\ 4,\ 6
\]
(for example, for $q=6$ only $p=1,5$, and for $q=9$ the values $p=1,2,4,5,7,8$). The total is $\boxed{27}$ different values. (For example $\dfrac{1}{2}=\dfrac{2}{4}=\dfrac{3}{6}=\dfrac{4}{8}$ are four fractions with one value.)
